\chapter{Using the Program}
\label{sec:usage}

\section{Building}

The program is built by Maven, which also runs the unit tests and creates a runnable archive:

\begin{verbatim}
mvn package
java -jar target/program.jar
\end{verbatim}

\section{Command Line}

The program solves every benchmark file found in the input directory and writes the results into a subdirectory of the output directory for each run. The nine mandatory arguments of the original edition are kept (Table~\ref{tab:args}); they can be followed by the optional arguments of Table~\ref{tab:options}. Without the optional arguments, the program runs POEMS with the corrections of Section~\ref{sec:defects}; the modifications evaluated in Chapter~\ref{sec:testing} are switched on by the system properties of Table~\ref{tab:switches}.

\begin{verbatim}
program.jar IN OUT VERBOSE R I G N S BESTFIT [OPTIONS]
program.jar benchmark/ benchmark/result/ false 1 1000 500 3 50 true
program.jar benchmark/ benchmark/result/ false 5 0 0 2 0 true \
    --algorithm=lapoems --time=60
\end{verbatim}

\begin{table}
\begin{center}
\begin{tabularx}{\textwidth}{@{}lXl@{}}
\toprule
Name & Description & Example \\
\midrule
{\tt IN} & input directory & benchmark/ \\
{\tt OUT} & output directory & benchmark/result/ \\
{\tt VERBOSE} & verbose mode & false \\
{\tt R} & number of runs for each benchmark & 1 \\
{\tt I} & number of iterations & 1000 \\
{\tt G} & number of generations & 500 \\
{\tt N} & niche count (length of action sequences; LA-POEMS: maximal number of actions in a step) & 3 \\
{\tt S} & niche size (number of individuals in a niche) & 50 \\
{\tt BESTFIT} & best-fit prototype (otherwise complete tree) & true \\
\bottomrule
\end{tabularx}
\end{center}
\caption{Mandatory command line arguments}
\label{tab:args}
\end{table}

\begin{table}
\begin{center}
\begin{tabularx}{\textwidth}{@{}lX@{}}
\toprule
Option & Description \\
\midrule
{\tt --algorithm=poems|lapoems} & the solver: POEMS (default) or late acceptance POEMS \\
{\tt --time=SECONDS} & time limit of each run; without it, LA-POEMS makes $I \cdot G$ steps \\
{\tt --seed=NUMBER} & seed of the random number generator \\
{\tt --history=LENGTH} & LA-POEMS: late acceptance with this history length (0 = automatic) instead of simulated annealing \\
\bottomrule
\end{tabularx}
\end{center}
\caption{Optional command line arguments}
\label{tab:options}
\end{table}

For each run, the output directory contains the final floorplan as an SVG image ({\tt floorplan.svg}, and {\tt floorplan\_s.svg} with the tree drawn over it), the tree as a DOT graph, its evaluation, the prototype, and the floorplans of all improvements of POEMS, which can be assembled into an animation.

\section{Benchmark Files}

A benchmark file is a plain text file that starts with the number of modules, followed by the width and the height of every module (the format of the benchmarks distributed with the floorplanning tools of the University of Michigan \cite{bench}). The dimensions may be written as real numbers, but their fractional parts must be zero. The whole MCNC hp benchmark is shown in Fig.~\ref{fig:hpcode}.

\begin{figure}
\begin{minipage}{\textwidth}
\begin{verbatim}
11
1036.00 462.00
378.00 700.00
980.00 210.00
980.00 210.00
980.00 210.00
3304.00 546.00
3304.00 546.00
2016.00 252.00
3080.00 462.00
2016.00 252.00
3080.00 462.00
\end{verbatim}
\end{minipage}
\caption{The MCNC hp benchmark file}
\label{fig:hpcode}
\end{figure}

\section{Reproducing the Experiments}

The benchmark runner solves one benchmark with a given seed and time limit (in milliseconds) and prints one line: the benchmark, the seed, the dead space in percent, the unused, total and used area, the width and height, the number of evaluations and iterations, the time, and the result of the validation of the floorplan.

\begin{verbatim}
java -cp target/program.jar program.Benchmark \
    benchmark/n100.txt SEED TIME_MS I G N S
java -cp target/program.jar -Dpoems.algorithm=lapoems \
    program.Benchmark benchmark/n100.txt SEED TIME_MS 0 0 N L
\end{verbatim}

For LA-POEMS, {\tt N} is the maximal number of actions in a step and {\tt L} the late acceptance history length, which is used only with {\tt -Dpoems.t0=0} (0 = automatic).

The settings of the improved POEMS and of LA-POEMS that were evaluated in Chapter~\ref{sec:testing} are switched by Java system properties, listed in Table~\ref{tab:switches}. The directory {\tt experiments} of the repository \cite{repository} contains the scripts that ran all experiments of Chapter~\ref{sec:testing}, the raw results of every run, and the scripts that create the tables and figures of this thesis from them.

\begin{table}
\begin{center}
\small
\begin{tabularx}{\textwidth}{@{}llX@{}}
\toprule
Property & Default & Meaning \\
\midrule
{\tt poems.algorithm} & {\tt poems} & {\tt poems} or {\tt lapoems} \\
{\tt poems.fast} & true & POEMS: array evaluation (identical results, faster) \\
{\tt poems.dedupe} & false & POEMS: discard children whose sequence is in the population \\
{\tt poems.patience} & 0 & POEMS: end an iteration after $P$ generations without improvement (0 = never) \\
{\tt poems.first} & false & POEMS: end an iteration when the prototype is improved \\
{\tt poems.lahc} & 0 & POEMS: late acceptance of prototypes with a history of $L$ iterations (0 = off) \\
{\tt poems.move} & false & POEMS: use the move node action \\
{\tt poems.t0}, {\tt poems.t1} & $10^{-2}$, $10^{-5}$ & LA-POEMS: initial and final temperature of simulated annealing, relative to the area of the modules ($T_0 = 0$: late acceptance) \\
{\tt poems.ratio} & 20000 & LA-POEMS: automatic history length = expected steps / $C$ \\
{\tt poems.restart} & 0 & LA-POEMS: restart after $R \cdot n$ steps without improvement (0 = never) \\
{\tt poems.local} & 0.5 & LA-POEMS: probability of a local second module \\
{\tt poems.weights} & uniform & LA-POEMS: relative probabilities of the seven action types \\
{\tt poems.adaptive} & false & LA-POEMS: adaptive selection of the action types \\
{\tt poems.trace} & -- & benchmark runner: write the convergence trace into a file or directory \\
{\tt poems.placement} & -- & benchmark runner: write the final floorplan into a directory \\
{\tt poems.evaluations} & 0 & benchmark runner, LA-POEMS: stop after this number of steps (0 = time limit only) \\
\bottomrule
\end{tabularx}
\end{center}
\caption{System properties of the experimental settings}
\label{tab:switches}
\end{table}
